% =============================================================================
% KI-Hilfsmittelverzeichnis
% -----------------------------------------------------------------------------
% PFLICHT, sobald KI-Werkzeuge bei der Erstellung der Arbeit genutzt wurden
% (FOM-Vorgaben zu KI in Studium und Lehre). Hier wird OFFENGELEGT, welche
% Werkzeuge wofür eingesetzt wurden. Diese Offenlegung ersetzt KEINE
% Quellenangabe: KI-gestützte Aussagen müssen zusätzlich mit Fachliteratur
% belegt werden.
%
% Trage jedes tatsächlich genutzte Werkzeug mit Anbieter, Version und
% Einsatzzweck ein. Die konkreten Prompts gehören (in Kurzform) in den Anhang.
% =============================================================================
\section*{KI-Hilfsmittelverzeichnis}
\addcontentsline{toc}{section}{KI-Hilfsmittelverzeichnis}

\begin{itemize}
  \item \textbf{[Werkzeug, Version (Anbieter)]:} [Einsatzzweck, z.\,B.
        Strukturierung der Gliederung, Recherche und Gegenprüfung von Quellen,
        sprachliche Überarbeitung einzelner Absätze].
  \item \textbf{[Werkzeug, Version (Anbieter)]:} [Einsatzzweck].
  \item \textbf{[Werkzeug, Version (Anbieter)]:} [Einsatzzweck].
\end{itemize}

Die durch KI erzeugten Ergebnisse wurden als Arbeitsgrundlage genutzt und
anschließend eigenständig geprüft, bewertet und überarbeitet. Die Verantwortung
für Inhalt, Quellenauswahl und Schlussfolgerungen liegt beim Verfasser.

% -----------------------------------------------------------------------------
% Hinweis: Prüfe in deinem aktuellen FOM-Leitfaden bzw. den Modulvorgaben die
% exakt geforderte Form (manche Standorte verlangen eine tabellarische
% Darstellung mit Spalten: Werkzeug | Zweck | betroffene Kapitel | Datum).
% -----------------------------------------------------------------------------
